% ============================================================================
% Fichier cible 1 : ARIEL_LOGMINER_MEMOIRE_FINAL/chapters/chapitre2_etat_art.tex
%   (§1.8.2 « CICIDS2017 », PDF p. 22)
% Ligne à modifier : 336 (« mesurées~\cite{engelen2021troubleshooting,
%   lanvin2023errors}. La chute du F1 observée ici entre séparation aléatoire
%   et séparation par scénario (section~6.2) ... »).
% Fichier cible 2 : ARIEL_LOGMINER_MEMOIRE_FINAL/chapters/chapitre5_resultats.tex
%   Insérer un label après la ligne 53 (« \section{CICIDS2017 : influence du
%   protocole} »), la section n'en ayant pas.
% Auteur : Agent 1 (intégrité factuelle), chantier 2026-11-05.
% Justification : la section 6.2 du PDF est « Valeur de la spécialisation »
%   (main.toc) ; la chute du F1 CICIDS2017 est présentée en section 4.2
%   (« CICIDS2017 : influence du protocole », PDF p. 81). Un renvoi par
%   \ref évite qu'une renumérotation recrée l'erreur.
% ============================================================================

% --- (1) chapitre5_resultats.tex, après la ligne 53 : ajouter
\label{sec:cicids-resultats}

% --- (2) chapitre2_etat_art.tex, ligne 336 : remplacer la ligne par
mesurées~\cite{engelen2021troubleshooting,lanvin2023errors}. La chute du F1 observée ici entre séparation aléatoire et séparation par scénario (section~\ref{sec:cicids-resultats}) est donc cohérente avec ces constats indépendants portant sur le même corpus.
